\documentclass[UTF8,a4paper,12pt]{ctexart}

% 独立编译版本仅供附录版式复核；正式提交使用已经合并附录的 main.tex。
\usepackage[top=2.5cm,bottom=2.5cm,left=2.5cm,right=2.5cm]{geometry}
\usepackage{amsmath,amssymb}
\usepackage{booktabs,longtable,array}
\usepackage{xurl}
\usepackage[hidelinks]{hyperref}
\usepackage{xcolor}
\usepackage{fvextra}
\usepackage{multicol}
\usepackage{fancyhdr}
\usepackage{titlesec}

\newcommand{\TeamControlNumber}{TJMML2026393}
\pagestyle{fancy}
\fancyhf{}
\setlength{\headheight}{15pt}
\fancyhead[R]{\small 队伍控制号：\TeamControlNumber}
\fancyfoot[C]{\thepage}
\renewcommand{\headrulewidth}{0pt}
\renewcommand{\footrulewidth}{0pt}

% 正文当前结束于第25页；正文页数变化后只需修改此处。
\setcounter{page}{26}
\setlength{\parindent}{2em}
\setlength{\parskip}{0pt}
\linespread{1.25}
\Urlmuskip=0mu plus 1mu

\renewcommand{\thesection}{\Alph{section}}
\renewcommand{\thesubsection}{\Alph{section}.\arabic{subsection}}
\titleformat{\section}{\centering\bfseries\fontsize{15pt}{20pt}\selectfont}
  {附录\thesection\quad}{0pt}{}

\newcommand{\CodeFile}[2]{%
  \par\medskip
  \noindent\textbf{#1}\par
  \noindent\textbf{路径：}\nolinkurl{#2}\par\smallskip
  \VerbatimInput[
    fontsize=\tiny,
    numbers=left,
    numbersep=3pt,
    frame=lines,
    framesep=1mm,
    rulecolor=\color{black!30},
    breaklines=true,
    breakanywhere=true,
    tabsize=4
  ]{#2}
}

\begin{document}
\appendix

\section{支撑材料文件列表}

支撑材料以 ZIP 压缩包单独提交。为便于评阅和复核，表~\ref{tab:file-list}仅列出文件类别、核心入口及用途；完整目录结构、运行命令、依赖版本和随机参数见压缩包根目录的 \texttt{README.md}。赛题提供的原始附件不重复列入自主数据资料。

\small
\begin{longtable}{
  >{\centering\arraybackslash}p{0.06\textwidth}
  >{\raggedright\arraybackslash}p{0.40\textwidth}
  >{\raggedright\arraybackslash}p{0.43\textwidth}
}
  \caption{支撑材料文件列表}\label{tab:file-list}\\
  \toprule
  序号 & 文件或目录 & 内容及用途 \\
  \midrule
  \endfirsthead
  \multicolumn{3}{c}{表\thetable（续）}\\
  \toprule
  序号 & 文件或目录 & 内容及用途 \\
  \midrule
  \endhead
  \bottomrule
  \endlastfoot

  1 & \nolinkurl{README.md}、各问题 \nolinkurl{README.md} & 总体目录、运行环境、运行顺序、输入输出与结果复现说明。\\
  2 & \nolinkurl{requirements.txt}、\nolinkurl{environment.yml}、\nolinkurl{config.json} & Python依赖、Conda环境及问题一搜索参数。\\

  3 & \nolinkurl{question1/program/run.py} & 问题一统一运行入口：机队规模搜索、路线优化、结果输出与复核。\\
  4 & \nolinkurl{question1/program/src/uav_q1/} & 问题一启发式求解、精确模型、独立校验及Excel输出模块；完整代码见附录B。\\
  5 & \nolinkurl{question1/results/} & 问题一正式工作簿、结构化路线和搜索记录。\\

  6 & \nolinkurl{question2/program/q2_solver.py} & 问题二负载均衡求解与独立校验程序；完整代码见附录C。\\
  7 & \nolinkurl{question2/results/} & 问题二工作簿、结构化方案、汇总指标和逐架路线。\\

  8 & \nolinkurl{question3/program/q3_solver.py} & 问题三算法A：直飞、必要等待与路线优化。\\
  9 & \nolinkurl{question3/program/q3_solver_enhanced.py} & 问题三算法B：时变可见图、空间绕飞与两阶段优化。\\
  10 & \nolinkurl{question3/program/q3_fast_validator.py}、\nolinkurl{q3_enhanced_validator.py} & 算法A快速校验与算法B逐事件独立复核；完整代码见附录D。\\
  11 & \nolinkurl{question3/results/} & 算法A对照结果、算法B正式结果及机队规模敏感性结果。\\

  12 & \nolinkurl{generated_figures/}、各问题辅助绘图与工作簿构建程序 & 论文图件、结果整理和格式转换文件，仅作为电子支撑材料提交。\\
  13 & \nolinkurl{AI_tool_usage_details.pdf} & 人工智能工具名称、用途、人工修改和核验情况。\\
\end{longtable}
\normalsize

\noindent\textbf{结果口径：}问题三正式方案采用算法B，固定问题一最终采用的机队规模 $N_0=(4,2,5,4)$；算法A仅作基准对照，$N_0+1$、$N_0+2$ 仅用于敏感性分析。全部正式结果均由独立程序复核。验证通过仅说明方案满足本文实现的约束，不代表已经证明全局最优。

\section{问题一源程序}

以下五个文件共同构成问题一完整程序，运行入口为 \texttt{run.py}。测试脚本、环境安装脚本、说明文档和预计算结果仅放入电子支撑材料。

\begin{multicols*}{2}
\CodeFile{问题一统一运行入口}{submission_package/question1/program/run.py}
\CodeFile{问题一启发式求解核心模块}{submission_package/question1/program/src/uav_q1/solver_engine.py}
\CodeFile{问题一固定机队规模精确模型}{submission_package/question1/program/src/uav_q1/exact_milp.py}
\CodeFile{问题一独立结果校验模块}{submission_package/question1/program/src/uav_q1/validation.py}
\CodeFile{问题一工作簿输出模块}{submission_package/question1/program/src/uav_q1/excel_output.py}
\end{multicols*}

\section{问题二源程序}

问题二求解、结果复算及CSV输出均集中在同一程序中；用于生成论文图件和最终Excel版式的辅助脚本仅放入电子支撑材料。

\begin{multicols*}{2}
\CodeFile{问题二负载均衡求解与校验程序}{submission_package/question2/program/q2_solver.py}
\end{multicols*}

\section{问题三源程序}

以下收入算法A、算法B及两套独立验证程序。算法B直接调用算法A文件中的数据结构和基础时空冲突函数，因此二者共同构成完整可运行程序。绘图、结果筛选和工作簿格式转换脚本仅放入电子支撑材料。

\begin{multicols*}{2}
\CodeFile{问题三算法A主程序}{submission_package/question3/program/q3_solver.py}
\CodeFile{问题三算法B主程序}{submission_package/question3/program/q3_solver_enhanced.py}
\CodeFile{问题三快速独立校验程序}{submission_package/question3/program/q3_fast_validator.py}
\CodeFile{问题三算法B逐事件独立验证程序}{submission_package/question3/program/q3_enhanced_validator.py}
\end{multicols*}

\end{document}
